% ID: SEL-D0007-BC-E05
% ESTADO: validada. S008, 2026-11-07; tres pasadas del mismo autor y revisión matemática, didáctica y visual.
\textbf{a) Partición y probabilidad total.}\par\nopagebreak[4]
Sea $S$ ``generar anticuerpos'' y $P$ ``recibir placebo''. Los grupos forman una partición de los 5000 voluntarios:
\[
P(A)=3000/5000=0.6,\quad P(B)=1500/5000=0.3,\quad P(P)=500/5000=0.1.
\]
Usando las frecuencias del ejercicio, $P(S\mid A)=0.90$, $P(S\mid B)=0.95$ y $P(S\mid P)=0$. Así,
\[
P(S)=0.6(0.90)+0.3(0.95)+0.1(0)=0.54+0.285=\boxed{0.825}.
\]

\textbf{b) Placebo sabiendo ausencia de anticuerpos.}\par\nopagebreak[4]
$P(S^c)=0.175$ y todos los que recibieron placebo están en $S^c$. Aplicando Bayes,
\[
P(P\mid S^c)=\frac{P(P)P(S^c\mid P)}{P(S^c)}
=\frac{0.1(1)}{0.175}=\boxed{\frac47\approx0.571429}.
\]
Comprobación por recuentos: generan anticuerpos 2700 de A y 1425 de B, total 4125. No los generan 300 de A, 75 de B y 500 de placebo, total 875; $500/875=4/7$. Son resultados del conjunto descrito en el problema, sin extrapolar los porcentajes a otros contextos.
